% =============================================================================
% Chapter 12: ML Operations at Scale — Vol II Lecture Slides
% =============================================================================
% IMAGES:
%   - n-models-complexity.pdf    - platform-roi.pdf
%   - mlops-maturity.pdf         - model-registry-deps.pdf
%   - ml-cicd-pipeline.pdf       - staged-rollout.pdf
%   - monitoring-pyramid.pdf     - feature-freshness.pdf
% =============================================================================
\documentclass[aspectratio=169, 12pt]{beamer}
\usepackage{../../assets/beamerthememlsys}

\mlsyssetup{
  volume       = {Volume II},
  chapter      = {Chapter 12},
  logo         = {../../assets/img/logo-mlsysbook.png},
  instlogo     = {../../assets/img/logo-harvard.png},
  chaptertitle = {ML Operations at Scale},
}

% --- Fonts ---
\usepackage[T1]{fontenc}
\usepackage[scaled=0.9]{helvet}
\usepackage{courier}
\renewcommand{\familydefault}{\sfdefault}

% --- Packages ---
\usepackage{booktabs}
\usepackage{amsmath}

% --- Image paths ---
\graphicspath{
  {images/}
}

% --- Chapter-specific macros ---
\newcommand{\TCO}{\text{TCO}_{\text{ML}}}

% --- Helper: safe image include ---
\newcommand{\safeimg}[2][width=\textwidth,keepaspectratio]{%
  \IfFileExists{images/#2}{\includegraphics[#1]{#2}}{%
    \IfFileExists{#2}{\includegraphics[#1]{#2}}{%
      \fbox{\parbox[c][2.5cm][c]{0.85\linewidth}{%
        \centering\footnotesize\textcolor{midgray}{[Missing image]}}}%
    }%
  }%
}

% --- Section count ---
\setcounter{mlsystotalsections}{6}

\title{ML Operations at Scale}
\author{Vijay Janapa Reddi}
\institute{Harvard University}
\date{}

\begin{document}

% =============================================================================
% TITLE SLIDE
% =============================================================================
\mlsystitle{ML Operations at Scale}{From Single-Model to Platform Engineering}{cover_ops_scale.png}

% =============================================================================
% LEARNING OBJECTIVES
% =============================================================================
\begin{frame}{Learning Objectives}
\note{
% -- LINK: Connect to prior chapters
Students built serving infrastructure in Part III. This chapter asks:
what happens when you manage not one model, but a hundred?

% -- NARRATE: What to SAY
Read each objective aloud, pausing on ``platform ROI,'' ``feature store
architectures'' (LO 7 --- new this year), and ``TCO framework.'' These three
anchor the quantitative and architectural reasoning for the entire chapter.

% -- ENGAGE: Specific question
Ask: ``How many of you have deployed more than one model to production?
At what count did ad hoc practices start breaking?''
Give 10 seconds for a show of hands.

% -- WARN: Specific misconception
Students assume operations scale linearly with model count.
Correct: dependencies grow as O(N^2), alerts as O(N*M).

% -- FLEX: [CORE]
[CORE] Never skip --- objectives frame the entire lecture.
IF AHEAD: Ask students to rank which objective they most want to master.
IF SHORT: Read objectives without elaboration, move on.
}

\small
\begin{enumerate}
  \item Calculate \textbf{platform ROI} to justify shared ML infrastructure
  \item Design \textbf{dependency-aware model registries} for multi-model portfolios
  \item Implement \textbf{CI/CD pipelines} tailored to model risk profiles
  \item Architect \textbf{hierarchical monitoring} that prevents alert fatigue
  \item Quantify \textbf{ML technical debt} using deployment velocity and toil metrics
  \item Compare \textbf{centralized}, \textbf{embedded}, and \textbf{hybrid} organizational patterns
  \item Evaluate \textbf{feature store architectures} for freshness SLOs and point-in-time correctness
  \item Apply the \textbf{TCO framework} to make infrastructure investment decisions
\end{enumerate}

\end{frame}

\begin{frame}{Visual Language}
\note{
% -- NARRATE: What to SAY
Point to each card: ``Blue = compute, green = data, orange = routing,
red = error. These are consistent across every diagram in this course.
When you see red in a pipeline diagram, something is bottlenecked.''

% -- FLEX: [OPTIONAL]
[OPTIONAL] Skip if students have seen this in a previous chapter deck.
IF SHORT: Say ``same color system as last lecture'' and advance.
}

\small
Throughout this course, colors carry meaning:

\vspace{0.3cm}
\begin{columns}[T]
  \begin{column}{0.45\textwidth}
    \begin{mlsyscard}{computestroke}
    \textbf{Blue} --- Compute / Processing\\
    {\footnotesize GPU ops, forward/backward pass, inference}
    \end{mlsyscard}
    \vspace{0.15cm}
    \begin{mlsyscard}{datastroke}
    \textbf{Green} --- Data / Memory\\
    {\footnotesize Data flow, caches, healthy paths}
    \end{mlsyscard}
  \end{column}
  \begin{column}{0.45\textwidth}
    \begin{mlsyscard}{routingstroke}
    \textbf{Orange} --- Routing / Scheduling\\
    {\footnotesize Load balancers, batch windows}
    \end{mlsyscard}
    \vspace{0.15cm}
    \begin{mlsyscard}{errorstroke}
    \textbf{Red} --- Error / Cost / Bottleneck\\
    {\footnotesize Loss, decode phase, waste}
    \end{mlsyscard}
  \end{column}
\end{columns}
\end{frame}



% =============================================================================
\section{Platform Operations}
% =============================================================================

\begin{frame}{The N-Models Problem}
\note{
% -- LINK: Connect to prior concept
Students just saw the learning objectives listing platform ROI and dependency
management. This slide makes the problem visceral: why platforms exist.

% -- NARRATE: What to SAY
Point to the diagram: ``At 10 models, a few shared data sources. At 50,
the dependency graph is a hairball. At 100, a single upstream change
cascades unpredictably.'' Trace the quadratic growth curve.

% -- ENGAGE: Specific question
Ask: ``At your organization, how many models share the same data sources?''
Expected answer: most students underestimate --- typical is 5--10 shared sources.

% -- WARN: Specific misconception
Students assume managing 100 models is 100x the work of managing one.
Correct: dependencies grow as O(N^2), so 100 models is closer to 10,000x
the coordination complexity.

% -- FLEX: [CORE]
[CORE] This motivates the entire chapter.
IF AHEAD: Ask ``At what N does your dependency graph become unmanageable?''
IF SHORT: Show diagram, state O(N^2), move on.
}

% --- Full-width diagram ---
\centering
\safeimg[width=0.92\textwidth,height=4.5cm,keepaspectratio]{n-models-complexity.pdf}

\vspace{0.1cm}
\mlsysinsight{Key:}{Per-model practices collapse at $\sim$50 models. Dependencies grow as $O(N^2)$.}

\end{frame}

\begin{frame}{Operational Complexity Growth}
\note{
% -- LINK: Connect to prior concept
The N-models diagram showed the complexity curve. This table puts
concrete operational labels on each step of that curve.

% -- NARRATE: What to SAY
Walk column by column: ``At 1 model, monitoring is a single dashboard.
At 100, you have 100 dashboards nobody reads. Debugging shifts from
local to distributed tracing --- a qualitatively different skill.''

% -- ENGAGE: Specific question
Ask: ``Which column transitions from manageable to unmanageable first?''
Expected answer: monitoring (it is the first to break because alert
volume grows as N times M metrics).

% -- WARN: Specific misconception
Students focus on deployment coordination but miss that monitoring
breaks first. Alert fatigue precedes deployment chaos.

% -- FLEX: [OPTIONAL]
[OPTIONAL] This table reinforces the N-models diagram.
IF SHORT: Point to the monitoring row, state the key insight, advance.
IF AHEAD: Ask students to fill in a ``1000 models'' column mentally.
}

\footnotesize
\renewcommand{\arraystretch}{1.15}
\begin{tabular}{@{}llll@{}}
  \toprule
  \textbf{Aspect} & \textbf{1 Model} & \textbf{10 Models} & \textbf{100 Models} \\
  \midrule
  \textbf{Deployment}  & None             & Ad hoc            & Critical path \\
  \textbf{Dependencies}& None             & Some overlap      & Dense graph \\
  \textbf{Monitoring}  & 1 dashboard      & 10 dashboards     & Unmanageable \\
  \textbf{On-call}     & Single team      & Multi-team        & Org-wide \\
  \textbf{Utilization} & Often idle       & Moderate sharing  & Efficiency critical \\
  \textbf{Debugging}   & Local            & Cross-team        & Distributed tracing \\
  \bottomrule
\end{tabular}

\vspace{0.15cm}
\mlsysalert{Consequence:}{100 models with independent failure modes create a constant stream of alerts that masks genuine emergencies.}

\end{frame}

\begin{frame}{Quantifying Platform ROI}
\note{
% -- LINK: Connect to prior concept
The complexity table showed operations becoming unmanageable. This slide
answers: what is the economic case for investing in a platform?

% -- NARRATE: What to SAY
Point to the equation: ``N is model count, T-saved is hours per model,
C-eng is engineer cost. The numerator grows linearly with N; the
denominator is fixed.'' Walk through the worked example: ``50 models,
40 hours each, \$150/hr = \$3.6M before, \$1.6M after. 56\% savings.''

% -- ENGAGE: Specific question
Ask: ``At what model count does your organization need a platform team?''
Give 10 seconds. Expected: most say 50--100; the surprise is 20--50.

% -- WARN: Specific misconception
Students think platform ROI is linear. Correct: it is superlinear because
shared infrastructure amortizes over all models simultaneously.

% -- FLEX: [CORE]
[CORE] The ROI equation is the quantitative anchor for the chapter.
IF AHEAD: Ask ``What if T-saved is only 10 hours instead of 30?''
IF SHORT: Show equation, state the 56\% number, skip the worked example detail.
}

\small
\begin{columns}[T]
  \begin{column}{0.48\textwidth}
    \safeimg[width=\textwidth,height=5.5cm,keepaspectratio]{platform-roi.pdf}
  \end{column}
  \begin{column}{0.48\textwidth}
    $$\text{ROI}_{\text{platform}} = \frac{N \times T_{saved} \times C_{eng}}{C_{platform}}$$

    \vspace{0.15cm}
    {\footnotesize
    \textbf{Worked example} (50 models):
    \begin{itemize}\setlength\itemsep{0pt}
      \item Before: 50 $\times$ 40 hrs $\times$ \$150 = \$3.6M/yr
      \item After: 50 $\times$ 10 hrs $\times$ \$150 + \$667K = \$1.6M/yr
      \item \textcolor{datastroke}{\textbf{Savings: 56\%}} (\$2M/yr)
    \end{itemize}
    }

    \vspace{0.1cm}
    \mlsysconcept{Threshold:}{Break-even at 20--50 models.}
  \end{column}
\end{columns}

\end{frame}

% --- ACTIVE LEARNING: Micro-Retrieval Cue ---
\begin{frame}{Quick Check}
\note{
% -- NARRATE: What to SAY
Pause 15 seconds. Then cold-call. Answer: 20--50 models.
Say: ``Most students guess 100+. The surprise is the threshold is
much lower because shared infrastructure amortizes across all models.''

% -- FLEX: [OPTIONAL]
[OPTIONAL] Micro-retrieval cue reinforcing the ROI slide.
IF SHORT: Skip entirely --- the predict exercise covers this ground.
}

\centering
\vspace{1.0cm}
{\Large\bfseries Quick Check}

\vspace{0.6cm}
{\large At what model count does platform investment\\typically break even?}

\vspace{0.5cm}
{\normalsize\textcolor{midgray}{15 seconds --- then I will cold-call.}}

\end{frame}



\begin{frame}{MLOps Maturity Hierarchy}
\note{
% -- LINK: Connect to prior concept
The ROI equation showed that platforms pay for themselves. This slide
asks: what does the maturity journey look like?

% -- NARRATE: What to SAY
Point to each level: ``L0 = scripts on laptops. L1 = per-model CI/CD,
the most common state. L2 = shared platform, where the superlinear
returns kick in. L3 = enterprise governance across the org.'' Emphasize
the L1-to-L2 transition: ``This is where most organizations stall.''

% -- ENGAGE: Specific question
Ask: ``Where is your organization on this staircase?''
Show of hands for each level. Most will cluster at L0--L1.

% -- WARN: Specific misconception
Students think the jump from L0 to L1 is the hard part. Correct:
L0-to-L1 is just adding CI/CD. The L1-to-L2 jump requires
organizational change --- shared infrastructure, common APIs, platform team.

% -- FLEX: [OPTIONAL]
[OPTIONAL] Reinforces the platform investment argument.
IF SHORT: Show staircase, name the four levels, emphasize L1-to-L2, advance.
}

% --- Full-width diagram ---
\centering
\safeimg[width=0.88\textwidth,height=4.5cm,keepaspectratio]{mlops-maturity.pdf}

\vspace{0.1cm}
\small
\begin{columns}[T]
  \begin{column}{0.22\textwidth}
    \centering
    \textcolor{errorstroke}{\textbf{L0: Manual}}\\
    {\scriptsize Ad hoc, 1--2 models}
  \end{column}
  \begin{column}{0.22\textwidth}
    \centering
    \textcolor{routingstroke}{\textbf{L1: Per-Model}}\\
    {\scriptsize CI/CD per model}
  \end{column}
  \begin{column}{0.22\textwidth}
    \centering
    \textcolor{computestroke}{\textbf{L2: Platform}}\\
    {\scriptsize Shared infra}
  \end{column}
  \begin{column}{0.22\textwidth}
    \centering
    \textcolor{datastroke}{\textbf{L3: Enterprise}}\\
    {\scriptsize Org-wide governance}
  \end{column}
\end{columns}

\end{frame}

% --- ACTIVE LEARNING 1: Predict ---
\begin{frame}{Predict: When Does Platform Investment Pay Off?}
\note{[2 min] Prediction exercise. Students write for 60 seconds.
Do NOT reveal the answer yet. This primes the platform economics discussion.
Ask 2--3 students to share.}

\centering
\vspace{1.0cm}
{\Large\bfseries Think--Write--Share}

\vspace{0.6cm}
{\large Your organization has 30 models, each costing\\
40 engineer-hours/month in operational toil.\\[0.3cm]
\alert{At what model count does a \$2M platform pay for itself?}}

\vspace{0.6cm}
{\normalsize Write your estimate. \textcolor{midgray}{(60 seconds)}}

\end{frame}

% =============================================================================
\section{Multi-Model Management}
% =============================================================================

\begin{frame}{Dependency-Aware Model Registry}
\note{
% -- LINK: What prior concept connects to this slide
The N-models problem showed dependencies grow as O(N-squared). This slide
makes dependency cascades concrete: one silent upstream change poisons
all downstream consumers.

% -- NARRATE: What to SAY while showing this slide
[3 min] Walk through the dependency graph: ``When the embedding model
updates silently, every downstream model receives garbage features.
Without a dependency-aware registry, organizations discover these
cascades through production failures --- not proactive monitoring.''

% -- ENGAGE: Specific question, prediction, or task for THIS slide
Ask: ``What happens when the embedding model updates silently?''
Give 10 seconds. [Expected: garbage predictions for all downstream models.]

% -- WARN: What students will get wrong on THIS topic
Common error: students think pinning model versions prevents cascades.
Correct framing: pinning delays the problem but creates technical debt ---
eventually you must propagate updates through the entire graph.

% -- FLEX: [CORE] + contingency
[CORE] Dependency management is the central multi-model challenge.
IF AHEAD: Ask ``How would you test a model update against all its consumers?''
IF SHORT: Show diagram, state the cascade failure mode, move on.
}

% --- Full-width diagram ---
\centering
\safeimg[width=0.88\textwidth,height=4.5cm,keepaspectratio]{model-registry-deps.pdf}

\vspace{0.1cm}
\mlsysalert{Without tracking:}{Organizations discover dependencies through production failures.}

\end{frame}

\begin{frame}{Ensemble Management}
\note{
% -- LINK: Connect to prior concept
The Dependency-Aware Model Registry showed that cascades happen when the
embedding model updates silently. Ensemble management is the operational
answer: how do you safely update one component of a 10--50 model pipeline?

% -- NARRATE: What to SAY
[3 min] Walk through the deployment stages: ``Shadow logs but does not
serve---this is where you verify the new component does not produce
garbage outputs. Canary routes 1\% of live traffic. Staged ramps from
5\% to 100\% while monitoring business metrics. Soak watches for
delayed effects: a personalization model may look fine at 100\% traffic
for three days but start degrading on day seven as feedback loops kick in.''
Point to the interaction effects callout: ``A component update that
improves its local metric---say, the retrieval model improves recall by
5\%---might degrade the full system because the ranking model was trained
on the old retrieval distribution. The ranking model has never seen
the new set of candidates.''

% -- ENGAGE: Specific question
Ask: ``What happens if the ranking model updates but the retrieval model
does not?'' Give 15 seconds. Expected: the ranking model now scores
candidates from a distribution it was never trained on, potentially
producing low-quality or unsafe rankings despite passing all unit tests.

% -- WARN: Specific misconception
Students assume updating one model in an ensemble is safe as long as
individual model metrics improve. Correct: the ensemble is a joint
distribution. Changing one component changes the joint, and the other
components were trained on the old joint. System-level metrics
(engagement, revenue, safety) are the only reliable signal.

% -- FLEX: [OPTIONAL]
[OPTIONAL] Ensemble interaction effects are important for RecSys heavy courses.
IF SHORT: Show staged deployment table; state the interaction-effects warning; advance.
IF AHEAD: Ask ``How would you use shadow mode to safely test a new retrieval model?''
}

\footnotesize
\textbf{A single recommendation request invokes 10--50 models:}

\vspace{0.1cm}
\renewcommand{\arraystretch}{1.1}
\begin{tabular}{@{}llrl@{}}
  \toprule
  \textbf{Deploy Stage} & \textbf{Actions} & \textbf{Duration} & \textbf{Rollback Trigger} \\
  \midrule
  \textbf{Shadow}       & Log only, not served        & 24--48 h  & Quality below threshold \\
  \textbf{Canary}       & 1\% traffic                 & 4--8 h    & Statistical regression \\
  \textbf{Staged}       & 5\% $\to$ 25\% $\to$ 100\% & 24--72 h  & Business metric drop \\
  \textbf{Soak}         & Full traffic, extended watch & 7--14 d   & Delayed effects \\
  \bottomrule
\end{tabular}

\vspace{0.15cm}
\begin{mlsyscard}{routingstroke}
{\footnotesize\textbf{Interaction effects:} A component change improving local metrics might degrade system-level performance through subtle interactions.}
\end{mlsyscard}

\end{frame}

\begin{frame}{Model-Type Operations Diversity}
\note{[2 min] Different model types need fundamentally different operational
patterns. Fraud detection needs hourly updates; LLMs need monthly staged
rollouts. The principle: risk profile determines operational cadence.}

\footnotesize
\renewcommand{\arraystretch}{1.15}
\begin{tabular}{@{}lllll@{}}
  \toprule
  \textbf{Model Type} & \textbf{Update Freq.} & \textbf{Deploy} & \textbf{Primary Risk} & \textbf{Rollback} \\
  \midrule
  \textbf{LLM}   & Monthly     & Staged, careful  & Quality, safety     & Hours--days \\
  \textbf{RecSys} & Daily--weekly & Shadow, A/B     & Engagement drop     & Minutes \\
  \textbf{Fraud}  & Hourly--daily & Rapid + rollback & False negatives     & Seconds \\
  \textbf{Vision} & Weekly--monthly & Canary         & Accuracy regression & Minutes \\
  \bottomrule
\end{tabular}

\vspace{0.15cm}
\mlsysconcept{Principle:}{Risk profile determines operational cadence.}

\vspace{0.1cm}
{\footnotesize LLMs operate slowly because quality regressions are hard to detect. Fraud detection operates continuously because adversaries do not wait.}

\end{frame}

% =============================================================================
\section{CI/CD at Scale}
% =============================================================================

\begin{frame}{ML CI/CD Pipeline}
\note{[3 min] Walk through the pipeline stages and gates. Key difference
from software CI/CD: must validate data, model performance, and fairness ---
not just code correctness. Each stage is idempotent. Ask: ``Which gate
is most likely to catch a silent regression?''}

% --- Full-width diagram ---
\centering
\safeimg[width=0.92\textwidth,height=4.8cm,keepaspectratio]{ml-cicd-pipeline.pdf}

\vspace{0.1cm}
\mlsysinsight{Key difference:}{ML CI/CD validates \textbf{data + model + fairness}, not just code.}

\end{frame}

\begin{frame}{Staged Rollout Strategy}
\note{
% -- LINK: Connect to prior concept
The ML CI/CD pipeline showed the stages that gate a model before promotion.
This slide shows what happens after the gate opens: how to release safely
when unit tests cannot guarantee real-world behavior.

% -- NARRATE: What to SAY
[3 min] Walk through the rollout diagram: ``Shadow logs but does not serve.
Canary routes 1\% of live traffic. Staged ramps from 5\% to 25\% to 100\%.
Soak watches for delayed effects at full traffic over 7--14 days.''
Point to the canary duration formula: ``At 1M req/hr with a 1\% canary,
you need 10,000 samples to detect a statistically significant regression,
which takes a minimum of 1 hour. The total confident rollout is about
4 hours. An instant blue-green deployment skips all of this.''
State the 20$\times$ risk reduction: ``A 5\% canary limits blast radius
by 20$\times$ relative to full rollout. If the model is broken, only 5\%
of users see the failure, not 100\%.''

% -- ENGAGE: Specific question
Ask: ``What does the 20$\times$ risk reduction protect against?'' Give 10
seconds. Expected: it limits the blast radius of a regression -- it does
NOT prevent the regression from occurring, it only controls how many users
are affected before detection.

% -- WARN: Specific misconception
Students assume canary deployment prevents regressions. Correct: canary
only limits blast radius -- it reduces the percentage of users exposed to
a bad model before rollback, but it does not make detection faster.
The regression still occurs; canary just contains it. Fast detection
requires good monitoring (alert system from the previous section).

% -- FLEX: [CORE]
[CORE] Staged rollouts are a quantitative safety argument, not just best practice.
IF SHORT: Show diagram, state 20$\times$ blast radius reduction; skip formula.
IF AHEAD: Ask ``What if your 1\% canary is demographically unrepresentative of the full user base?''
}

% --- Full-width diagram ---
\centering
\safeimg[width=0.88\textwidth,height=4.0cm,keepaspectratio]{staged-rollout.pdf}

\vspace{0.1cm}
\small
\textbf{Canary duration:} $t_{stage} = \frac{n_{samples}}{r_{requests} \times p_{stage}}$

\vspace{0.1cm}
{\footnotesize At 1M req/hr with 1\% canary: 10K samples needed $\to$ 1 hour minimum.\\
Total confident rollout: $\sim$4 hours (vs.\ instant blue-green with no safety net).}

\end{frame}

% --- ACTIVE LEARNING 2: Exercise ---
\begin{frame}{Your Turn: Cost of a Delayed Alert}
\note{[3 min] Give students 90 seconds. Answer: 432M requests $\times$ 0.005
lost CTR $\times$ \$0.50/click = \$1.08M. The insight: a ``minor'' 0.5\%
regression at scale is a million-dollar daily disaster.}

\small
\begin{columns}[T]
  \begin{column}{0.55\textwidth}
    {\normalsize\bfseries Calculate the Loss}

    \vspace{0.2cm}
    A recommendation model has a silent bug reducing CTR by 0.5\% (5.0\% $\to$ 4.5\%).
    \begin{itemize}
      \item 5,000 QPS, each click worth \$0.50
      \item Detection takes 24 hours
    \end{itemize}

    \vspace{0.15cm}
    \textbf{How much revenue is lost?}

    {\footnotesize\textcolor{midgray}{(90 seconds --- then compare with a neighbor)}}
  \end{column}
  \begin{column}{0.42\textwidth}
    \pause
    \begin{mlsyscard}{errorstroke}
    \textbf{Solution:}\\[0.1cm]
    {\footnotesize
    Requests: 5K $\times$ 86,400 = 432M\\
    Lost clicks: 432M $\times$ 0.005 = 2.16M\\
    Loss: 2.16M $\times$ \$0.50\\[0.1cm]
    $=$ \textcolor{errorstroke}{\textbf{\$1,080,000}}\\[0.1cm]
    Every hour of faster detection saves \textbf{\$45K}.
    }
    \end{mlsyscard}
  \end{column}
\end{columns}

\end{frame}

\begin{frame}{CI/CD Patterns by Model Type}
\note{[2 min] Connect back to model-type diversity. LLMs: quality-gated,
days to weeks. Fraud: threshold-gated, hours with seconds rollback.
The key: one-size-fits-all CI/CD fails at platform scale.}

\footnotesize
\renewcommand{\arraystretch}{1.15}
\begin{tabular}{@{}lllll@{}}
  \toprule
  \textbf{Pattern} & \textbf{Model Type} & \textbf{Validation} & \textbf{Rollout} & \textbf{Rollback} \\
  \midrule
  \textbf{Quality-gated}   & LLM   & Human eval + safety  & Days--weeks & Hours \\
  \textbf{Metric-driven}   & RecSys & Engagement metrics  & Hours--days & Minutes \\
  \textbf{Threshold-gated} & Fraud  & Precision/recall    & Hours       & Seconds \\
  \textbf{Accuracy-focused} & Vision & Classification     & Days        & Minutes \\
  \bottomrule
\end{tabular}

\vspace{0.15cm}
\begin{mlsyscard}{crimson}
{\footnotesize \textbf{Rollout risk:} $R = P_{\text{regression}} \times I_{\text{regression}} \times E_{\text{exposure}}$. Reduce any factor to reduce total risk.}
\end{mlsyscard}

\end{frame}

% =============================================================================
\section{Monitoring at Scale}
% =============================================================================

\mlsysfocus{The False Alarm Tax}{%
$P(\text{false alert}) = 1 - (1-\alpha)^N$\\[0.3cm]
{\normalsize With $\alpha=0.05$, $N=1000$: $P \approx 1.0$}%
}

\begin{frame}{Alert Fatigue: The Math}
\note{[3 min] This is the key quantitative insight. Even at 3-sigma
specificity, 100 models with 10 metrics generate 864 false alerts/day.
At 2-sigma: 14,400/day. Operators learn to ignore everything.
Ask: ``How many alerts per day before your on-call team gives up?''}

\small
\textbf{100 models $\times$ 10 metrics $\times$ 5-min checks:}

\vspace{0.15cm}
{\footnotesize
\renewcommand{\arraystretch}{1.2}
\begin{tabular}{@{}lrrr@{}}
  \toprule
  \textbf{Threshold} & \textbf{FPR} & \textbf{Daily False Alerts} & \textbf{Status} \\
  \midrule
  2-sigma (95\%)   & 5.0\%  & 14,400 & \textcolor{errorstroke}{\textbf{Useless}} \\
  3-sigma (99.7\%) & 0.3\%  & 864    & \textcolor{routingstroke}{\textbf{Overwhelming}} \\
  4-sigma (99.99\%) & 0.01\% & 29    & \textcolor{datastroke}{\textbf{Manageable}} \\
  \bottomrule
\end{tabular}
}

\vspace{0.15cm}
\mlsysalert{Destructive dynamic:}{Operators ignore alerts $\to$ genuine issues get lost $\to$ monitoring provides \emph{negative} value.}

\end{frame}

\begin{frame}{Hierarchical Monitoring Architecture}
\note{[3 min] The solution to alert fatigue. Four levels: business,
portfolio, model, infrastructure. High-level metrics trigger alarms;
lower-level metrics serve investigation only. Walk through the pyramid.}

\small
\begin{columns}[T]
  \begin{column}{0.48\textwidth}
    \safeimg[width=\textwidth,height=5.5cm,keepaspectratio]{monitoring-pyramid.pdf}
  \end{column}
  \begin{column}{0.48\textwidth}
    \textbf{Four abstraction levels:}

    \vspace{0.15cm}
    \textcolor{errorstroke}{\textbf{L1 Business:}} Revenue, engagement\\
    {\scriptsize Few metrics, high signal, exec attention}

    \vspace{0.1cm}
    \textcolor{routingstroke}{\textbf{L2 Portfolio:}} RecSys lift, fraud rate\\
    {\scriptsize Domain-level aggregation}

    \vspace{0.1cm}
    \textcolor{computestroke}{\textbf{L3 Model:}} Accuracy, latency, drift\\
    {\scriptsize Investigation on anomaly only}

    \vspace{0.1cm}
    \textcolor{datastroke}{\textbf{L4 Infra:}} GPU util, feature store\\
    {\scriptsize Platform team routing}

    \vspace{0.1cm}
    \mlsysconcept{Principle:}{Alert upward, investigate downward.}
  \end{column}
\end{columns}

\end{frame}

\begin{frame}{The Control Plane That Disappeared (2021)}
\note{
% -- LINK: Connect to prior concept
The hierarchical monitoring pyramid showed how to detect problems. This war
story shows what happens when the control plane that runs your monitoring
shares the same failure domain as the systems it watches.

% -- NARRATE: What to SAY
[2 min] Walk through the four fields: ``On October 4, 2021, a single
misconfigured BGP command disconnected Facebook's data centers from each
other. The routing withdrawal also took down the DNS servers and internal
management tools needed for remote recovery. Engineers could not use
standard SSH or web consoles to fix the outage because those tools depended
on the same network plane that had just failed. Teams restored connectivity
over hours using constrained out-of-band paths.'' Point to the Systems
lesson: ``This is not just a Facebook problem. Your ML platform has the
same trap: training schedulers, model registries, feature stores, and
inference routers all share the same auth, DNS, and network plane. When
that plane fails, so does your ability to diagnose and recover.''

% -- ENGAGE: Specific question
Ask: ``If your model registry becomes unreachable during an incident, can
your on-call engineer still roll back a bad model without it? Do you have
runbooks that work when your tools are down?''
Give 15 seconds. Expected: most students have not thought through
out-of-band recovery for ML platform failures.

% -- WARN: Specific misconception
Students assume the control plane and data plane fail independently.
Correct: in shared cloud infrastructure, they often share DNS, auth, and
network fabric. An incident that takes down one can take down both.
Recovery tooling must be physically isolated to survive the same failure
mode it is designed to fix.

% -- FLEX: [CORE]
[CORE] War stories anchor abstract monitoring principles in real cost.
IF SHORT: Read Context and Systems lesson only; skip Consequence details.
IF AHEAD: Ask ``What would you isolate first to prevent this in your ML platform?''
}

\footnotesize
\begin{columns}[T]
  \begin{column}{0.48\textwidth}
    \textbf{Context:} Meta's global services depended on a backbone network
    carrying internal control traffic for Facebook, Instagram, and WhatsApp.

    \vspace{0.15cm}
    \textbf{Failure:} Oct 4, 2021 --- a BGP command intended to assess
    backbone capacity unintentionally disconnected data centers from each
    other. The routing withdrawal also broke DNS and internal recovery tools.

    \vspace{0.15cm}
    \textbf{Consequence:} Major services unavailable globally for hours;
    recovery required constrained out-of-band paths because standard
    management tools depended on the failed network plane.
  \end{column}
  \begin{column}{0.48\textwidth}
    \begin{mlsyscard}{errorstroke}
    {\footnotesize\textbf{Systems lesson:}\\[0.1cm]
    The control plane is a dependency with its own blast radius. ML
    platforms expose the same trap: training schedulers, model registries,
    feature stores, and inference routers share the auth, DNS, and network
    plane that an incident can take down.\\[0.1cm]
    \textbf{Required:} explicit isolation between the control plane and
    the ML systems that depend on it.}
    \end{mlsyscard}
  \end{column}
\end{columns}

\end{frame}

\begin{frame}{Drift Detection: PSI}
\note{[2 min] Two types of drift: covariate (input $P(X)$ changes) and
concept ($P(Y|X)$ changes). PSI is cheap to compute and serves as a leading
indicator. Thresholds: $<$0.1 stable, 0.1--0.25 investigate, $\geq$0.25 act.
Monitoring lag: 2\% accuracy drop needs $\sim$2,500 labeled samples.}

\footnotesize
\textbf{Population Stability Index} detects feature distribution shift:

\vspace{0.05cm}
$$PSI = \sum_{i=1}^{n} (A_i - E_i) \times \ln\left(\frac{A_i}{E_i}\right)$$

\vspace{0.05cm}
\renewcommand{\arraystretch}{1.1}
\begin{tabular}{@{}lll@{}}
  \toprule
  \textbf{PSI Value} & \textbf{Interpretation} & \textbf{Action} \\
  \midrule
  $< 0.1$          & No significant shift   & Continue monitoring \\
  $0.1$ -- $0.25$  & Moderate shift          & Investigate \\
  $\geq 0.25$      & Significant shift       & \textcolor{errorstroke}{\textbf{Immediate action}} \\
  \bottomrule
\end{tabular}

\vspace{0.1cm}
\mlsysalert{Monitoring lag:}{Detecting a 2\% accuracy drop needs $\sim$2,500 samples. At 1K labels/hr $\to$ \textbf{2.5 hrs} broken. PSI on inputs is faster.}

\end{frame}

% --- ACTIVE LEARNING 3: Discussion ---
\begin{frame}{Discussion: Which Monitoring Level Fires First?}
\note{[3 min] Turn-and-talk. Students discuss in pairs for 90 seconds.
The answer depends on the failure mode: infrastructure failures show at L4 first,
while data drift shows at L3 then propagates to L2/L1.}

\centering
\vspace{0.5cm}
{\large\bfseries Turn and Talk \textcolor{midgray}{(90 seconds)}}

\vspace{0.45cm}
{\large An upstream feature pipeline silently starts\\
returning null values for 3\% of users.\\[0.25cm]
\alert{Which monitoring level detects this first?\\
How long until the business metric dashboard shows a dip?}}

\vspace{0.45cm}
\begin{columns}[c]
  \begin{column}{0.22\textwidth}\centering\small L1\\Business\end{column}
  \begin{column}{0.22\textwidth}\centering\small L2\\Portfolio\end{column}
  \begin{column}{0.22\textwidth}\centering\small L3\\Model\end{column}
  \begin{column}{0.22\textwidth}\centering\small L4\\Infra\end{column}
\end{columns}

\end{frame}

% =============================================================================
\section{Features \& Organization}
% =============================================================================

\begin{frame}{Feature Freshness: Batch vs.\ Streaming}
\note{
% -- LINK: Connect to prior concept
The monitoring section showed how to detect when model behavior drifts.
Feature freshness is one of the fastest ways to cause that drift without
changing the model at all.

% -- NARRATE: What to SAY
[3 min] Point to the diagram: ``Freshness latency $L_{\text{freshness}}$
is the gap between when an event happens and when the model can use a
feature derived from it. Batch: 12--24 hours. Streaming: 1--5 seconds.
That is a 10,000$\times$ improvement, which translates to 10--20\%
engagement lift in recommendation systems because the model can react
to what users just did rather than what they did yesterday.'' Then pivot
to fraud: ``For fraud detection the same gap gives adversaries a 24-hour
window to exploit patterns the model does not yet see.''

% -- ENGAGE: Specific question
Ask: ``For fraud detection, how stale can features be before adversaries
exploit the window?'' Give 15 seconds. Expected: any lag beyond a few
minutes is exploitable for high-velocity card fraud.

% -- WARN: Specific misconception
Students assume streaming is always better and always worth the cost.
Correct: streaming is expensive to build, operate, and maintain. Batch
is cheaper and adequate for the majority of features (user demographics,
long-term preferences, slowly-changing signals). The right question is
not ``should we stream?'' but ``which features need freshness SLOs under
5 minutes, and for which can we tolerate 12 hours?'' Most features fall
in the batch category; streaming the wrong ones wastes engineering budget.

% -- FLEX: [OPTIONAL]
[OPTIONAL] Freshness motivates the dual-store architecture on the next slide.
IF SHORT: State the $L_{\text{freshness}}$ formula and the 24-hour fraud window; advance.
}

% --- Full-width diagram ---
\centering
\safeimg[width=0.88\textwidth,height=4.2cm,keepaspectratio]{feature-freshness.pdf}

\vspace{0.05cm}
\small
$L_{\text{freshness}} = T_{\text{available}} - T_{\text{event}}$

\vspace{0.05cm}
{\footnotesize Streaming: \textbf{10--20\% engagement lift}. For fraud detection, day-old features give adversaries a 24-hour exploitation window.}

\end{frame}

\begin{frame}{Feature Store: Dual-Store Architecture (1/2)}
\note{
% -- LINK: Connect to prior concept
The Feature Freshness slide showed the latency gap between batch and
streaming. This slide explains the architectural pattern that solves it:
why you need two separate stores, not one clever database.

% -- NARRATE: What to SAY
[2 min] Walk through the architecture: ``Training needs high-throughput
batch scans over years of history. Serving needs sub-10-ms point lookups
for individual requests. No single database efficiently satisfies both
constraints. The dual-store pattern resolves this: an offline store
(BigQuery, Snowflake, data lake) handles the terabyte-scale analytical
scans for training. An online store (Redis, DynamoDB, Bigtable) handles
the millisecond key-value lookups for inference. A materialization layer
keeps them in sync.''

% -- ENGAGE: Specific question
Ask: ``Why can't one database serve both training and inference?''
Give 10 seconds. Expected: training needs sequential scans over TB,
serving needs single-key lookups under 10 ms --- no DB optimizes both.

% -- WARN: Specific misconception
Common error: students assume streaming features into a single DB is
``simpler'' than two stores. Correct: a single DB suffers from impedance
mismatch -- either too slow for serving or too expensive for training.

% -- FLEX: [CORE]
[CORE] This slide covers LO 7 (feature store architectures). Do not skip.
IF SHORT: Focus on the invariant $f_{\text{train}}(x) \equiv f_{\text{serve}}(x)$.
IF AHEAD: Ask ``What is the role of the materialization layer?''
}

\footnotesize
\textbf{The core conflict:}
\begin{itemize}\setlength\itemsep{0pt}
  \item Training: TB-scale sequential scans, hours latency OK
  \item Serving: single-entity lookups, $<$10 ms p99
  \item One DB cannot efficiently satisfy both
\end{itemize}

\vspace{0.08cm}
\textbf{Dual-store pattern:}
\begin{itemize}\setlength\itemsep{0pt}
  \item \textcolor{datastroke}{\textbf{Offline store}} (BigQuery/S3/Iceberg): training data, high throughput
  \item \textcolor{computestroke}{\textbf{Online store}} (Redis/DynamoDB): serving, low latency
  \item \textbf{Materialization layer}: keeps both stores in sync
\end{itemize}

\vspace{0.08cm}
\mlsysconcept{Invariant:}{$f_{\text{train}}(x) \equiv f_{\text{serve}}(x)$ ---
{\footnotesize training-serving skew silently destroys accuracy when this breaks.}}

\end{frame}

\begin{frame}{Feature Store: Point-in-Time Correctness (2/2)}
\note{
% -- LINK: Connect to prior concept
The dual-store architecture ensures training and serving use the same
feature pipeline. This slide covers the subtlest correctness requirement:
point-in-time joins. Even with the right architecture, naive time handling
causes data leakage.

% -- NARRATE: What to SAY
[1 min] Explain point-in-time correctness: ``Training must join features
as they existed \emph{at the event timestamp}, not current values. A user
had 4 clicks at noon; a naive join supplies the end-of-day value of 10 ---
2.5$\times$ inflation and guaranteed production failure. At 10,000 models
sharing the same feature platform, data leakage risk multiplies. Point-in-time
joins must be an infrastructure guarantee, not a per-team check.''

% -- ENGAGE: Specific question
Ask: ``What happens to the fraud model if the online store's user
transaction feature is 10 minutes stale but training used 1-minute windows?''
Expected: the model scores stale features --- adversaries exploit the window.

% -- FLEX: [CORE]
[CORE] Point-in-time correctness is the key correctness invariant for feature stores.
IF SHORT: State the 2.5$\times$ inflation example and advance.
}

\begin{mlsyscard}{routingstroke}
{\small\textbf{Point-in-time correctness:}\\[0.1cm]
Training must join features as they existed \emph{at the event timestamp},
not current values.

\vspace{0.1cm}
\textbf{Example}: A user had 4 clicks at noon; a naive join supplies the
end-of-day value of 10 --- \textbf{2.5$\times$ inflation} and guaranteed
production failure.

\vspace{0.1cm}
At 10,000 models sharing the same feature platform, data leakage risk
multiplies. Point-in-time joins are an \textbf{infrastructure guarantee},
not a per-team check.}
\end{mlsyscard}

\end{frame}

\begin{frame}{Organizational Patterns for ML Platforms}
\note{
% -- LINK: Connect to prior concept
The Feature Store architecture requires a shared platform team to maintain
consistency guarantees. This slide asks: what organizational structure
should own and operate that platform?

% -- NARRATE: What to SAY
[2 min] Walk column by column: ``Centralized gives you consistency but
creates a bottleneck --- every model team waits for platform approvals.
Embedded gives velocity but fragments your feature store, monitoring, and
CI/CD into per-team silos that can't talk to each other. Hybrid is the
practical answer for 50--100 models: a core platform team owns shared
infrastructure (the dual-store feature store, hierarchical monitoring,
CI/CD gates) while domain leads embedded in product teams handle
domain-specific needs.''

% -- ENGAGE: Specific question
Ask: ``At 75 models across 3 product lines, which pattern would you choose
and why?'' Give 15 seconds. Expected: hybrid is the right answer, because
centralized would bottleneck 75 model teams and embedded would produce
75 inconsistent feature pipelines.

% -- WARN: Specific misconception
Students assume centralized is always better because consistency is good.
Correct: centralized works well above 100 models where the platform team
is fully staffed and specialized. Below that threshold, a two-person
platform team is a constant bottleneck for 75 model teams trying to ship.

% -- FLEX: [OPTIONAL]
[OPTIONAL] Organizational patterns are contextual.
IF SHORT: Show the table; state the 50--100 model hybrid threshold; advance.
IF AHEAD: Discuss what happens to organizational patterns when a startup
acquires a large ML platform --- how do embedded engineers integrate?
}

\footnotesize
\renewcommand{\arraystretch}{1.15}
\begin{tabular}{@{}lp{3.2cm}p{3.2cm}p{3.2cm}@{}}
  \toprule
  & \textbf{Centralized} & \textbf{Embedded} & \textbf{Hybrid} \\
  \midrule
  \textbf{Structure}  & Single platform team & ML eng in each team & Core + domain leads \\
  \textbf{Strength}   & Consistency          & Responsiveness       & Balanced \\
  \textbf{Risk}       & Bottleneck           & Fragmentation        & Coordination cost \\
  \textbf{Best for}   & 100+ models          & 10--20 models        & 50--100 models \\
  \bottomrule
\end{tabular}

\vspace{0.15cm}
\begin{mlsyscard}{computestroke}
{\footnotesize \textbf{Decision factors:} Model count, model similarity, org size, regulatory requirements, and infrastructure maturity determine the optimal pattern.}
\end{mlsyscard}

\end{frame}

\begin{frame}{Edge Fleet MLOps: Version Skew (1/2)}
\note{
% -- LINK: Connect to prior concept
The data-center MLOps patterns just covered assume controlled infrastructure.
This slide extends the same framework to millions of heterogeneous edge devices
where you cannot force an immediate rollout.

% -- NARRATE: What to SAY
[2 min] Walk through the two challenges: ``In a data center, a rollout
completes in hours. In an edge fleet, it takes weeks because devices
are offline, battery-saving, or on restricted networks. At any moment,
an organization may have 50 different model versions active. Every
upstream data pipeline must be backwards-compatible with models deployed
months ago.'' Then federated monitoring: ``You cannot stream raw predictions
back due to privacy and bandwidth --- devices compute local error rates and
transmit only aggregated, anonymized metrics. Fleet health must be reconstructed
from fragmented signals.''

% -- ENGAGE: Specific question
Ask: ``If you have 50 model versions active and want to add a new feature
to your upstream data pipeline, what constraints does version skew impose?''
Give 20 seconds. Expected: the new pipeline must produce outputs compatible
with all 50 active versions, or you must maintain parallel pipelines until
the oldest versions retire.

% -- WARN: Specific misconception
Students assume edge MLOps is just data-center MLOps with lower power.
Correct: the operational contract is fundamentally different. Recovery
is probabilistic, not deterministic.

% -- FLEX: [CORE]
[CORE] Covers Key Takeaways bullet on edge fleet management.
IF SHORT: State the 50-version-skew number; advance to HIL CI/CD slide.
IF AHEAD: Ask ``How would you design a canary rollout when 10\% of devices
are always offline?''
}

\footnotesize
\textbf{Challenge 1: Extreme version skew}
\begin{itemize}\setlength\itemsep{0pt}
  \item Data-center rollout: hours; edge fleet rollout: weeks to months
  \item Up to \textbf{50 concurrent model versions} active simultaneously
  \item Upstream data pipelines must stay backwards-compatible across all versions
\end{itemize}

\vspace{0.12cm}
\textbf{Challenge 3: Federated monitoring}
\begin{itemize}\setlength\itemsep{0pt}
  \item Cannot stream raw predictions back (privacy + bandwidth constraints)
  \item Devices compute local error rates; transmit only aggregated, anonymized metrics
  \item Fleet health must be reconstructed from fragmented signals
\end{itemize}

\vspace{0.1cm}
\mlsysalert{Key difference:}{Edge recovery is probabilistic, not deterministic --- a bad model takes weeks to roll back, not seconds.}

\end{frame}

\begin{frame}{Edge Fleet MLOps: HIL CI/CD (2/2)}
\note{
% -- LINK: Connect to prior concept
Version skew showed that edge rollouts take weeks. This slide covers the
quality gate that prevents bad models from entering that slow propagation path.

% -- NARRATE: What to SAY
[1 min] ``Passing accuracy gates is not enough for edge. A model must also
be validated on physical or emulated target hardware: specific NPU architectures,
microcontrollers, DSPs. A model that passes accuracy but exceeds the 1 MB memory
limit of a budget microcontroller, or triggers thermal throttling on a specific
SoC, is rejected at the HIL stage, not in production. Edge failures discovered
in production cannot be rolled back in seconds --- they take weeks to propagate.''

% -- FLEX: [CORE]
[CORE] HIL CI/CD is the key distinction between edge and data-center MLOps.
IF SHORT: State the three rejection criteria and the weeks-to-rollback cost.
}

\begin{mlsyscard}{computestroke}
{\footnotesize\textbf{Challenge 2: Hardware-in-the-Loop (HIL) CI/CD}\\[0.05cm]
Accuracy gates are \textbf{necessary but not sufficient} for edge deployments.\\[0.05cm]
CI/CD must include a \textbf{HIL stage}: validate on physical or emulated fleet hardware
\emph{before} promotion.\\[0.05cm]
A model is \textbf{rejected} if it:
\begin{itemize}\setlength\itemsep{0pt}
  \item Exceeds 1 MB memory on a budget microcontroller
  \item Triggers thermal throttling on a specific SoC
  \item Exceeds latency budget on a target NPU/DSP
\end{itemize}
Edge failures discovered in production cannot be rolled back in seconds ---
they take \textbf{weeks to propagate} through the fleet.}
\end{mlsyscard}

\end{frame}

\begin{frame}{TCO Framework for ML Systems}
\note{[3 min] The four-component cost structure evolves with scale. Early-stage:
iteration-dominated (optimize velocity). Production: inference-dominated
(optimize efficiency). Walk through the equation and the shift.
Ask: ``Which cost component dominates at your organization?''}

\small
$$\TCO = C_{\text{train}} + C_{\text{infer}} + C_{\text{data}} + C_{\text{iter}}$$

\vspace{0.15cm}
{\footnotesize
\renewcommand{\arraystretch}{1.2}
\begin{tabular}{@{}lll@{}}
  \toprule
  \textbf{Component} & \textbf{Early Stage} & \textbf{Production Scale} \\
  \midrule
  $C_{\text{train}}$  & Dominant (experimentation)     & Fixed (scheduled retraining) \\
  $C_{\text{infer}}$  & Negligible (few users)         & \textcolor{errorstroke}{\textbf{Dominant}} (70--80\% of TCO) \\
  $C_{\text{data}}$   & Small (curated datasets)       & Growing (storage + pipelines) \\
  $C_{\text{iter}}$   & \textcolor{routingstroke}{\textbf{Dominant}} (engineer time)  & Reduced by automation \\
  \bottomrule
\end{tabular}
}

\vspace{0.15cm}
\mlsysconcept{Strategy:}{Match optimization investment to current cost structure.}

\end{frame}

\begin{frame}{Quantifying ML Technical Debt}
\note{
% -- LINK: What prior concept connects to this slide
The TCO framework showed cost structure. Technical debt is the hidden cost
that erodes deployment velocity --- Sculley et al.\ called it ``the high-interest
credit card of ML.''

% -- NARRATE: What to SAY while showing this slide
[3 min] Walk through the debt categories: ``Configuration debt: 3--5x more config
code than model code. Data debt: unstable input features that break silently.
Pipeline debt: glue code connecting systems nobody owns. Monitoring debt:
\$750K/yr from 15 incidents at \$50K each.'' Point to the audit framework:
``Measure deployment velocity (deploys/week), toil ratio (manual/total hours),
and incident rate. Teams with high debt see 70--80\% capacity consumed by
maintenance.''

% -- ENGAGE: Specific question, prediction, or task for THIS slide
Ask: ``What fraction of your team's time goes to maintaining existing models
vs.\ building new ones?'' Give 10 seconds. [Expected: most say 50\%+. The
surprise is that 70--80\% is common at mature organizations without debt management.]

% -- WARN: What students will get wrong on THIS topic
Common error: students think technical debt is inevitable and should wait.
Correct framing: debt compounds --- monitoring debt alone costs \$750K/yr.
Early investment in debt reduction has superlinear ROI.

% -- FLEX: [CORE] + contingency
[CORE] Technical debt quantification is a key learning objective.
IF AHEAD: Ask ``Which debt category would you audit first?''
IF SHORT: Show the table and state the 70--80\% maintenance number.
}

\footnotesize
\textbf{ML Technical Debt Categories} (Sculley et al., NeurIPS 2015):

\vspace{0.1cm}
\renewcommand{\arraystretch}{1.1}
{\scriptsize
\begin{tabular}{@{}lllr@{}}
  \toprule
  \textbf{Debt Category} & \textbf{Symptom} & \textbf{Metric} & \textbf{Annual Cost} \\
  \midrule
  \textbf{Configuration} & 3--5$\times$ more config than model code & Config LOC ratio & Eng.\ velocity \\
  \textbf{Data}          & Unstable input features break silently  & Feature staleness & Data quality \\
  \textbf{Pipeline}      & Glue code nobody owns                  & Ownership gaps    & Incident rate \\
  \textbf{Monitoring}    & 15 incidents $\times$ \$50K each        & MTTR, false alerts & \textbf{\$750K/yr} \\
  \bottomrule
\end{tabular}
}

\vspace{0.1cm}
\begin{mlsyscard}{errorstroke}
{\footnotesize \textbf{Debt audit metrics:} Deployment velocity (deploys/week), toil ratio (manual/total hrs), incident rate. Teams with unmanaged debt: \textbf{70--80\% capacity on maintenance}.}
\end{mlsyscard}

\end{frame}

% --- ACTIVE LEARNING: Alert System Design Exercise ---
\begin{frame}{Your Turn: Design an Alert System}
\note{
% -- LINK: What prior concept connects to this slide
Alert fatigue showed 14,400 false alerts/day at 2-sigma. This exercise asks
students to calculate false alert rates at different thresholds for a real
portfolio size.

% -- NARRATE: What to SAY while showing this slide
[3 min] Read the parameters. Give 90 seconds. After pause reveal: ``At 3-sigma,
FPR = 0.003, checks = 200 x 8 x 288 = 460,800, false alerts = 1,382/day.
At 4-sigma, FPR = 0.0001, false alerts = 46/day --- manageable.''

% -- ENGAGE: Specific question, prediction, or task for THIS slide
Students must calculate false alert rates at two thresholds. Key insight:
the jump from 3-sigma to 4-sigma cuts false alerts by 30x.

% -- WARN: What students will get wrong on THIS topic
Common error: students apply the same threshold to all metrics. Correct
framing: business-critical metrics need 3-sigma; infrastructure metrics
can use 4-sigma with investigation-only routing.

% -- FLEX: [CORE] + contingency
[CORE] Quantitative exercise for the monitoring section.
IF SHORT: Reduce to 60 seconds; show the 3-sigma answer only.
}

\small
\begin{columns}[T]
  \begin{column}{0.55\textwidth}
    {\normalsize\bfseries Design the Alert Thresholds}

    \vspace{0.15cm}
    \textbf{Portfolio}: 200 models, 8 metrics each.\\
    \textbf{Check interval}: 5 minutes (288/day).

    \vspace{0.1cm}
    \textbf{Calculate daily false alerts at:}
    \begin{enumerate}
      \item 3-sigma (FPR = 0.003)
      \item 4-sigma (FPR = 0.0001)
    \end{enumerate}

    {\footnotesize\textcolor{midgray}{(90 seconds --- then compare with a neighbor)}}
  \end{column}
  \begin{column}{0.42\textwidth}
    \pause
    \begin{mlsyscard}{computestroke}
    \textbf{Solution:}\\[0.1cm]
    {\footnotesize
    Checks/day: $200 \times 8 \times 288$\\
    = 460,800 checks\\[0.1cm]
    3-sigma: $460{,}800 \times 0.003$\\
    = \textcolor{errorstroke}{\textbf{1,382 false alerts/day}}\\[0.1cm]
    4-sigma: $460{,}800 \times 0.0001$\\
    = \textcolor{datastroke}{\textbf{46 false alerts/day}}\\[0.1cm]
    30$\times$ reduction!
    }
    \end{mlsyscard}
  \end{column}
\end{columns}

\end{frame}

% =============================================================================
\section{Wrap-Up}
% =============================================================================

\begin{frame}{Fallacies}
\note{[2 min] Four common misconceptions with quantitative evidence.}

\footnotesize
\textbf{Fallacy:} \textit{Operational complexity grows linearly with model count.}\\
100 models $\times$ 10 metrics at 5\% FPR $\to$ 14,400 false alerts/day. Dependencies grow as $O(N^2)$.

\vspace{0.08cm}
\textbf{Fallacy:} \textit{Platform investment only makes sense after 100+ models.}\\
A \$2M platform breaks even at $\sim$20 models. By 100, migration cost is 3--5$\times$ the initial build.

\vspace{0.08cm}
\textbf{Fallacy:} \textit{Monitoring training metrics provides sufficient observability.}\\
20 ms upstream slowdown $\to$ 30\% latency degradation across 10--50 downstream models.

\vspace{0.08cm}
\textbf{Fallacy:} \textit{Technical debt is inevitable and should wait.}\\
Monitoring debt: \$750K/yr (15 incidents $\times$ \$50K). Teams see 70--80\% capacity on maintenance.

\end{frame}

\begin{frame}{Pitfalls}
\note{[2 min] Three operational pitfalls. If short: cover just the first two.}

\footnotesize
\textbf{Pitfall:} \textit{Applying independent alerting to every model and metric.}\\
With $N=1000$ independent tests at $\alpha=0.05$: $P(\text{false alert}) = 1 - 0.95^{1000} \approx 1.0$. Operators ignore the noise; genuine incidents vanish. Use \textbf{hierarchical monitoring}.

\vspace{0.15cm}
\textbf{Pitfall:} \textit{Treating all deployments with uniform rollout procedures.}\\
Fraud detection needs hourly updates with seconds-fast rollback. LLMs need multi-day shadow testing. \textbf{Risk-based policies} must match model type.

\vspace{0.15cm}
\textbf{Pitfall:} \textit{Defaulting to batch feature computation for all features.}\\
Daily batch yields $L_{\text{freshness}} \approx 12$--24 hrs; streaming achieves 1--5 seconds. For recommendations, streaming yields \textbf{10--20\% engagement lift}; for fraud, batch gives adversaries a 24-hour window.

\end{frame}

% --- RETRIEVAL PRACTICE ---

% --- MUDDIEST POINT ---
\begin{frame}{Muddiest Point}
\note{[2 min] Quick anonymous poll. Students write on a slip of paper or submit
digitally. Collect and scan for patterns. Address the top 2--3 confusions in the
next lecture's opening. This closes the feedback loop.}

\centering
\vspace{1.0cm}
{\Large\bfseries What was the \alert{muddiest point} today?}

\vspace{0.8cm}
{\normalsize Write down the concept you found \textbf{most confusing}.}

\vspace{0.5cm}
{\small\textcolor{midgray}{Anonymous. One sentence. Submit before you leave.}}

\end{frame}

\begin{frame}{What Were the Key Ideas?}
\note{[2 min] Retrieval practice. Students write 90 seconds, no notes.
Do NOT show next slide yet. Walk around the room.}

\centering
\vspace{1.5cm}
{\Large\bfseries Close your notes.}

\vspace{0.8cm}
{\large Write down the \textbf{4 most important concepts} from today.}

\vspace{0.8cm}
{\normalsize\textcolor{midgray}{90 seconds --- no peeking.}}

\end{frame}

\begin{frame}{Key Takeaways}
\note{[2 min] Reveal. Walk through each bullet. Emphasize quantitative
anchors: 56\% savings, 14,400 alerts, 20$\times$ risk reduction,
\$1.08M daily loss.}

\scriptsize
\begin{itemize}\setlength\itemsep{0pt}
  \item \textbf{Platform ROI is superlinear}: Break-even at 20--50 models. 50 models saves 56\% (\$2M/yr).
  \item \textbf{Dependencies grow as $O(N^2)$}: Per-model practices collapse at $\sim$50 models.
  \item \textbf{Risk profile determines cadence}: LLMs: months. Fraud: hours. One-size CI/CD fails.
  \item \textbf{Alert fatigue is mathematical}: 100 models $\times$ 10 metrics $\to$ 14,400 false alerts/day. Use hierarchical monitoring.
  \item \textbf{Staged rollouts reduce risk 20$\times$}: 5\% canary limits blast radius.
  \item \textbf{Feature stores enforce the invariant} $f_{\text{train}}(x) \equiv f_{\text{serve}}(x)$: dual offline/online architecture prevents training-serving skew at platform scale.
  \item \textbf{Edge fleet demands HIL CI/CD}: Up to 50 concurrent model versions; accuracy gates alone are insufficient --- validate on physical NPU/DSP before promotion.
  \item \textbf{TCO shifts with scale}: Early = iteration-dominated; production = inference-dominated.
\end{itemize}

\end{frame}

\begin{frame}{References}
\note{[1 min] Point students to canonical papers.}

\small
\mlsysref{Sculley+15}{Sculley et al. ``Hidden Technical Debt in ML Systems.'' NeurIPS 2015.}
\mlsysref{Covington+16}{Covington et al. ``Deep Neural Networks for YouTube Recommendations.'' RecSys 2016.}
\mlsysref{Chapelle+12}{Chapelle et al. ``Large-Scale Validation and Analysis of Interleaved Search Evaluation.'' TOIS 2012.}
\mlsysref{Kohavi+09}{Kohavi et al. ``Controlled Experiments on the Web.'' KDD 2009.}
\mlsysref{Hardt+16}{Hardt et al. ``Equality of Opportunity in Supervised Learning.'' NeurIPS 2016.}
\mlsysref{Paleyes+22}{Paleyes et al. ``Challenges in Deploying ML.'' ACM Comp.\ Surveys, 2022.}

\end{frame}

\begin{frame}{Next Lecture: Security and Privacy}
\note{[1 min] Forward hook. A globally distributed fleet that autonomously
adapts presents an expansive attack surface. Data poisoning, model extraction,
and differential privacy are the next challenges.}

\centering
\vspace{0.8cm}
{\large A fleet that is globally distributed\\
and autonomously adapting presents an\\
\alert{expansive attack surface.}}

\vspace{0.6cm}
\begin{columns}[c]
  \begin{column}{0.30\textwidth}
    \centering
    \begin{mlsyscard}{errorstroke}
    {\small\bfseries Data Poisoning}\\[0.1cm]
    {\scriptsize Corrupt training data}
    \end{mlsyscard}
  \end{column}
  \begin{column}{0.30\textwidth}
    \centering
    \begin{mlsyscard}{computestroke}
    {\small\bfseries Model Extraction}\\[0.1cm]
    {\scriptsize Steal model via API}
    \end{mlsyscard}
  \end{column}
  \begin{column}{0.30\textwidth}
    \centering
    \begin{mlsyscard}{datastroke}
    {\small\bfseries Differential Privacy}\\[0.1cm]
    {\scriptsize Protect training data}
    \end{mlsyscard}
  \end{column}
\end{columns}

\vspace{0.4cm}
{\normalsize\textbf{How do we defend the fleet?}}

\end{frame}



\appendix

\begin{frame}{Backup: Platform ROI Sensitivity Analysis}
\note{
% -- NARRATE: What to SAY while showing this slide
Backup. Shows how ROI changes with engineer cost and model count. Use when
students ask ``what if our engineers cost less?'' or ``we only have 15 models.''

% -- FLEX: [OPTIONAL] + contingency
[OPTIONAL] Only show if students question the 20--50 model break-even.
}

\footnotesize
\textbf{Platform break-even varies with engineer cost and model count:}

\vspace{0.15cm}
\renewcommand{\arraystretch}{1.15}
\begin{tabular}{@{}lrrr@{}}
  \toprule
  \textbf{Engineer Cost} & \textbf{10 Models} & \textbf{30 Models} & \textbf{50 Models} \\
  \midrule
  \$100/hr & ROI: 0.5$\times$ & ROI: 1.4$\times$ & ROI: 2.3$\times$ \\
  \$150/hr & ROI: 0.7$\times$ & ROI: \textbf{2.1$\times$} & ROI: \textbf{3.4$\times$} \\
  \$200/hr & ROI: 1.0$\times$ & ROI: 2.8$\times$ & ROI: 4.5$\times$ \\
  \bottomrule
\end{tabular}

\vspace{0.15cm}
\begin{mlsyscard}{crimson}
{\footnotesize At \$150/hr engineer cost, break-even is $\sim$15 models. The threshold is lower than most teams expect because shared infrastructure amortizes across all models simultaneously.}
\end{mlsyscard}

\end{frame}

\begin{frame}{Backup: Monitoring Tier Design Patterns}
\note{
% -- NARRATE: What to SAY while showing this slide
Backup. Detailed monitoring tier configuration for production deployments.
Use when students ask how to implement hierarchical monitoring.

% -- FLEX: [OPTIONAL] + contingency
[OPTIONAL] Implementation detail for students building monitoring systems.
}

\footnotesize
\renewcommand{\arraystretch}{1.15}
\begin{tabular}{@{}lllr@{}}
  \toprule
  \textbf{Tier} & \textbf{Metrics} & \textbf{Alert Target} & \textbf{Threshold} \\
  \midrule
  \textbf{L1 Business}    & Revenue, engagement      & Executive + on-call & 4-sigma \\
  \textbf{L2 Portfolio}   & Domain-level aggregation  & Domain lead         & 3.5-sigma \\
  \textbf{L3 Model}       & Per-model accuracy, drift & Investigation only  & 3-sigma \\
  \textbf{L4 Infra}       & GPU util, feature store   & Platform team       & 3-sigma \\
  \bottomrule
\end{tabular}

\vspace{0.15cm}
\textbf{Key pattern:} L1/L2 alert humans. L3/L4 feed dashboards for investigation \emph{after} higher-tier alerts fire. This prevents the 14,400 false alerts/day problem.

\end{frame}


\end{document}
